\begin{frame}
  \frametitle{Outline}
% One \nocite per line -- a multi-key \nocite spanning lines breaks bibtex.
\nocite{FPR13mar}
\nocite{GPWMFH13lim}
\nocite{CDJ18enh}
Order-driven markets are trading venues organised around \emph{limit orders}:
the participants themselves supply the prices,
and the venue only matches them against one another.

We describe the mechanism precisely:
what a limit order is,
how the matching algorithm clears it,
what the resulting object --- the limit order book --- looks like,
and what it costs to trade against it.
\end{frame}

\subsection{Introduction}

\begin{frame}
  \frametitle{Two market infrastructures}
Market infrastructures fall in two categories:
\emph{dealer markets} (also known as quote-driven markets)
and
\emph{limit order markets} (also known as order-driven markets).

\pause

EURONEXT, the Swiss, Helsinki, Canadian, Toronto, Australian, Hong Kong, Shenzhen and Tokyo
exchanges operate as purely order-driven markets;
NYSE, NASDAQ and LSE are order-driven with some legacy quote-driven mechanisms.

\pause

Order-driven markets are transparent:
every participant observes in real time all the others' commitments to trade,
organised in the so-called \emph{limit order book}.
\end{frame}

\subsection{Order-driven markets}

\begin{frame}
  \frametitle{The limit order}
A limit order is the fundamental action a participant can perform.
It is a 4-tuple
\[
(t,q,p,d),
\]
where $t$ is time, $q$ is size, $p$ is price,
and $d=1$ (buy) or $d=-1$ (sell) is direction.

\pause

The price $p$ is the \emph{worst} price the sender accepts:
the highest at which she buys, the lowest at which she sells.
By regulation $p$ is an integer multiple of the tick size $\tickSizeOfLOB$.

\pause

A participant may also \emph{cancel} a previously submitted order,
withdrawing her commitment to trade.
\end{frame}

\begin{frame}
  \frametitle{Matching}
The incoming order $(t,q,p,d)$ is matched against active orders
with timestamp $s<t$ and opposite direction $-d$.

\pause

The active order $(s,\rho,\pi,-d)$ \emph{matches} it if $\pi d \leq p d$;
then $\rho \wedge q$ shares transact at the price $\pi$,
and the two orders are replaced by
\[
(s,\rho - \rho\wedge q, \pi, -d),
\qquad
(t,q - \rho\wedge q, p, d).
\]

\pause

At least one of them has null size.
If the incoming order still has size left, the search continues;
whatever cannot be filled is recorded in the book and waits.
\end{frame}

\begin{frame}
  \frametitle{Price-time priority}
The search follows the ordering
\[
(s,\rho,\pi,-d)<(s\derivative,\rho\derivative,\pi\derivative,-d)
\quad \text{ iff } \quad
(d\pi<d\pi\derivative)\vee((\pi=\pi\derivative)\wedge(s<s\derivative)).
\]
Better price first; at equal price, earlier timestamp first.

\pause

This ordering is a \emph{total} order on each side of a trading epoch ---
antisymmetry holds because no two orders share a timestamp.

\pause

So each side is a queue,
and a position in it is an asset that cancelling forfeits.
\end{frame}

\begin{frame}
  \frametitle{The order queues}
\begin{defi}
\label{def.orderQueue}
Let $E$ be a trading epoch. The ask and bid order queues at time $t$ are
\[
\askOrderQueue_t := \lbrace (s,q,p,-1)\in E : \text{active at time } t \rbrace,
\]
\[
\bidOrderQueue_t := \lbrace (s,q,p,+1)\in E : \text{active at time } t \rbrace.
\]
\end{defi}

\pause

An order is \emph{active} (or outstanding) at time $t$ if it was submitted by then
and has been neither executed nor cancelled.
\end{frame}

\begin{frame}
  \frametitle{The book is a price grid}
A limit order book is a grid of equally spaced prices, of spacing $\tickSizeOfLOB$,
at which the active orders sit.

\pause

Buy offers organise on the left (the \emph{bid side}),
sell offers on the right (the \emph{ask side}) ---
not by decree, but because any overlap would already have been matched away.

\pause

The whole configuration at time $t$ is
\[
\Big(\bestAskPrice_t, \; \bestBidPrice_t, \;
\lbrace (\nthBestAskSize[i]_t, \nthBestBidSize[i]_t) : i=1,2,\dots \rbrace \Big),
\]
with $\nthBestAskPrice[i]_t = \bestAskPrice_t + (i-1)\tickSizeOfLOB$
and $\nthBestBidPrice[i]_t = \bestBidPrice_t - (i-1)\tickSizeOfLOB$.

\pause

$\nthBestAskSize[i]_t$ is the \emph{size} of that level, the shares queuing there.
\emph{Volume} is the shares transacted over an interval:
a level has a size at every instant, volume accumulates between two.
Volume is not a change in size --- a level also shrinks by cancellation.
\end{frame}

\begin{frame}
  \frametitle{Spread, mid-price, imbalance}
Three derived quantities assess an order book:
\[
\LOBspread_t = \abs{\bestAskPrice_t - \bestBidPrice_t},
\qquad
\midPrice_t = \frac{\bestAskPrice_t + \bestBidPrice_t}{2},
\]

\pause

and the $n$-levels queue imbalance
\[
\queueImbalance^{n}_t =
\frac{\sum_{i\leq n}\nthBestBidSize[i]_t - \sum_{i\leq n}\nthBestAskSize[i]_t}{\sum_{i\leq n}\nthBestBidSize[i]_t + \sum_{i\leq n}\nthBestAskSize[i]_t}.
\]

\pause

The queue imbalance, called volume imbalance in \citet{CDJ18enh},
is widely accepted as a reliable signal for the next mid-price move:
close to $-1$ the mid-price will likely decrease, close to $+1$ it will likely increase.
\end{frame}

\begin{frame}
  \frametitle{Every limit order hides a market order}
Processing the sell limit order $(t,q,p,-1)$ is \emph{equivalent} to processing the pair
\[
\big[\,(t,\marketOrderSize,0,-1)\,,\;(t,q-\marketOrderSize,p,-1)\,\big],
\]
where
\[
\marketOrderSize = \min\Big(q, \textstyle\sum_{i\geq 1} \nthBestBidSize[i]_{t-} \indicator{\nthBestBidPrice[i]_{t-}\geq p}\Big).
\]

\pause

The first component has price $0$ --- immediate execution, nothing queued.
That is the \emph{market order}.
The second is exactly the part that rests in the book.

\pause

A market order \emph{walks the book} when $\marketOrderSize > \nthBestBidSize[1]_{t-}$:
it eats past the best level, each unit worse than the last.
Trades are market orders with $\marketOrderSize > 0$.
\end{frame}

\begin{frame}
  \frametitle{The cost of immediacy}
Executing a signed volume $\volume$ at time $t$ by \emph{market order}:
\[
\Delta \cashAccount_t^{\text{market order}}
= -\midPrice_{t}\volume
\underbrace{-\frac{\LOBspread_t}{2}\abs{\volume}}_{\text{we pay the spread}}.
\]

\pause

By \emph{limit order} sent at $t-1$ and hit at $t$:
\[
\Delta \cashAccount_t^{\text{limit order}}
= -\midPrice_{t-1}\volume
\underbrace{+\frac{\LOBspread_{t-1}}{2}\abs{\volume}}_{\text{we earn the spread}}.
\]

\pause

Same sign flipped on the spread ---
but notice the stale time index $t-1$ on the second line.
\end{frame}

\begin{frame}
  \frametitle{Adverse selection}
That stale index is the whole story. In portfolio value,
\[
\Delta \wealth_t^{\text{market order}}
= \inventory_{t-1}\Delta \midPrice_t - \frac{\LOBspread_t}{2}\abs{\volume},
\]
\[
\Delta \wealth_t^{\text{limit order}}
= \inventory_{t-1}\Delta \midPrice_t
+ \volume\, \Delta \midPrice_t
+ \frac{\LOBspread_{t-1}}{2}\abs{\volume}.
\]

\pause

The limit order earns the spread but carries the extra term $\volume\,\Delta\midPrice_t$:
the price move between sending and being hit.
We do not control when we are hit --- faster or better informed participants do.

\pause

This is \emph{adverse selection}.
The choice between market and limit orders rests on whether the spread
adequately compensates for it.
\end{frame}
